\documentclass[french]{book}
\usepackage{teach}
\usepackage{verbatim}
\usepackage{amsmath,amsfonts,amssymb}
\usepackage{pgfplots}
\usepackage{mwe}
\usepackage{yhmath} 
\usepackage{pstricks,pst-plot,pst-text,pst-tree,pst-eucl}
\usepackage[french,lined,boxed]{algorithm2e}	% pour les algorithmes
%\ configuration package algorithm
\SetKw{Retour}{Afficher}
\SetKwBlock{Traitement}{Début traitement}{Fin traitement.}
\SetKw{Afficher}{Afficher}
\usepackage{listings}
\newcommand{\arc}[1]{\wideparen{#1}}
\RequirePackage{graphicx}
\RequirePackage{ifpdf}
 \ifpdf
   \DeclareGraphicsRule{*}{mps}{*}{}
 \fi

\newcommand{\sysd}[2]{%
       \left\{
       \begin{aligned}
          #1\\
          #2\\
       \end{aligned}
       \right.%
       }
  \usepackage{fancyhdr}

  \usepackage{amsmath}
  \usepackage{amssymb}
  \usepackage{amsfonts}
  \usepackage{amsthm}
  \usepackage{mathrsfs}

  \usepackage{array}
  \usepackage{multicol}
  \setlength{\multicolsep}{3pt}

  \usepackage{graphicx}
  \graphicspath{{Figures/}}
  \usepackage{pstricks,pst-plot,pst-text,pst-tree,pst-eucl}
  \usepackage[all]{xy}

  \usepackage{fancybox}
  \usepackage{color}

%  \usepackage[frenchb]{babel}
  \usepackage{hyperref}

\usepackage{subfig}
\pgfplotsset{compat=newest}
\usepackage[all]{xy}
%\usepackage{geometry}
%\geometry{top=1cm, bottom=1cm, left=2cm, right=2cm}
\usepackage[utf8]{inputenc}
\usepackage[french]{babel}
\redefinecolor{thm}{title}{bgcolor}{alizarin}
\redefinecolor{prop}{title}{bgcolor}{meatbrown}
\redefinecolor{defn}{title}{bgcolor}{olivine}
\definecolor{alizarin}{rgb}{0.82, 0.1, 0.26}
\definecolor{meatbrown}{rgb}{0.9, 0.72, 0.23}
\definecolor{olivine}{rgb}{0.6, 0.73, 0.45}
\usepackage{mathrsfs}
%%
\newcommand{\R}{\mathbb {R}}
%\newcommand{\C}{\mathds {C}}
\newcommand{\N}{\mathbb {N}}
\newcommand{\Z}{\mathbb {Z}}
\newcommand{\Q}{\mathbb {Q}}
\newcommand{\D}{\mathbb {D}}
%%
\newcommand{\se}{\geq}
\newcommand{\ie}{\leq}
\newcommand{\tm}{\times}
%\newcommand{\ell}{l}
%\newcommand{\ell'}{l'}
%%
\newcommand{\barre}[1]{\overline{#1}}
\newcommand{\ds}{\displaystyle}
\newcommand{\limn}{\lim_{n\rightarrow +\infty}}
\newcommand{\limo}[1][x]{\ds\lim_{#1\rightarrow 0}}
\newcommand{\limite}[2][x]{\ds\lim_{#1\rightarrow #2}}
\newcommand{\limpinf}[1][x]{\ds\lim_{#1\rightarrow +\infty}}
\newcommand{\limminf}[1][x]{\ds\lim_{#1\rightarrow -\infty}}
\newcommand{\limiteg}[2][x]{\ds\lim_{#1\xrightarrow{<}#2}}
\newcommand{\limited}[2][x]{\ds\lim_{#1\xrightarrow{>}#2}}
%%
\newcommand{\vect}[1]{\overrightarrow{#1}}
\newcommand{\cro}[1]{ \ds [#1]}
\newcommand{\oij}{(\textit{O},{\overrightarrow{i}},{\overrightarrow{j}})}
%
\newcommand{\cp}[2]{%
        \begin{pmatrix}
        #1\\
        #2
        \end{pmatrix}%
        }
%
\newcommand{\calb}{\mathscr{B}}
\newcommand{\calc}{\mathscr{C}}
\newcommand{\cald}{\mathscr{D}}
\newcommand{\cale}{\mathscr{E}}
\newcommand{\calf}{\mathscr{F}}
\newcommand{\calp}{\mathscr{P}}
\newcommand{\cals}{\mathscr{S}}
\newcommand{\calt}{\mathscr{T}}
%
\newcommand{\suite}[1][u]{\left( #1_{n} \right)_{n\in\N}}
\newcommand{\suitep}[1][u]{\left( #1_{n} \right)_{n\in\N^{*}}}
\usetikzlibrary{arrows}

\newcommand{\poly}[1][a]{#1_0+#1_1x+\cdots+#1_nx^n}
\usepackage{tabularx}
\usepackage{array}
\usepackage{pst-tree}
\usepackage{pifont}
%\usepackage[cp1252]{inputenc}


%\newcommand{\binom}[2]{C_{#2}^{#1}}
\newcommand{\C}[2]{C_{#1}^{#2}}
\newcommand{\V}{\overrightarrow}
\newcommand{\Rep}{(O;\V{i};\V{j})}
%\newcommand{\Int}[2]{[\, #1 ; #2 \,]}
%\newcommand{\Reel}{\mathfrak{Re}} 
%\newcommand{\Ima}{\mathfrak{Im}} 
%\renewcommand{\Abs}[1]{\left \lvert#1\right \rvert}
\newcommand{\Cnp}{\begin{pmatrix} n\\p \end{pmatrix}}
%\newcommand{\C}[2]{\begin{pmatrix} #1\\#2 \end{pmatrix}}
\newcommand{\Syst}[2]{\left\{\begin{array}{lcccc} #1 \\ #2 \end{array}\right.}
\renewcommand{\arraystretch}{1.2} 
\newcolumntype{C}[1]{>{\centering\arraybackslash}p{#1cm}}
\newcommand{\dx}{dx}

\begin{document}


\tableofcontents


\chapter{Lois de probabilités discrètes}

\section{Loi uniforme}

\begin{defn}
Soit $E=\{1;2;...;n\}$ où $n$ est un entier naturel non nul. On dit qu'une variable aléatoire $X$ suit la \emph{loi uniforme} de paramètre $n$ si pour tout $k\in \{1;2;...;n\}$, 
$$P(X=k)=\frac{1}{n}$$
\end{defn}

\begin{demo}
Il faut montrer que l'on obtient bien une loi de probabilité en posant $P(X=k)=\frac{1}{n}$.\\
On a  $P(X=k)\geq 0$\\
et $\sum_{k=0}{n}\frac{1}{n}=\frac{1}{n}+\frac{1}{n}+...\frac{1}{n}=n\times \frac{1}{n}=1$.\\
Donc on définit bien une loi de probabilité. 
\end{demo}

\begin{rem}
La loi uniforme apparaît lorsque l'on modélise une expérience aléatoire à l'aide d'une loi \emph{équirépartie}, c'est à dire lorsque l'on a affaire à une situation \emph{d'équiprobabilité}.
\end{rem}



\begin{prop}
\begin{itemize}
\item[$\bullet$] Soit $A$ un événement de $E$. Alors $P(X\in A)=\frac{\textrm{nombre d'éléments de }A}{n}$.
\item[$\bullet$] L'espérance est $E(X) = \frac{n+1}{2}$
\end{itemize}

\end{prop}


\begin{demo}
Le premier point est immédiat. Preuve de l'espérance :\\
$E(X)=\sum_{k=1}^{n}kP(X=k)=\sum_{k=1}^{n}k\frac{1}{n}=\frac{1}{n}(\sum_{k=0}^{n}k)=\frac{1}{n}\frac{n(n+1)}{2}=\frac{n+1}{2}$\\
ou écrit autrement :
$E(X)=1\times P(X=1)+2\times P(X=2)+\cdots +n\times P(X=n)=1\times \frac{1}{n}+2\times \frac{1}{n}+\cdots +n\times \frac{1}{n}=\frac{1}{n}(1+2+\cdots+n)=\frac{1}{n}\frac{n(n+1)}{2}=\frac{n+1}{2}$.

\end{demo}

\section{Loi de Bernoulli}



\begin{defn}
Soit $p$ un nombre réel tel que $p\in[0;1]$. Soit $X$ une variable aléatoire. On dit que $X$ suit une loi de Bernoulli de paramètre $p$ si :
\begin{itemize}
 \item[$\bullet$] $X$ prend pour seules valeurs 1 (\og{}succès\fg{}) et 0 (\og{}échec\fg{}) ;
\item[$\bullet$] $P(X=1)=p$ et $P(X=0)=1-p$.
\end{itemize}
\end{defn}


\begin{prop}
Si $X$ suit une loi de Bernoulli de paramètre $p$ :
\begin{itemize}
\item[$\bullet$] L'espérance est $E(X)=p$ ;
\item[$\bullet$] la variance est $V(X)=p(1-p)$.
\end{itemize}
\end{prop}

\begin{demo}
\begin{itemize}
\item[$\bullet$] $E(X)=0\times P(X=0)+1\times P(X=1)=1\times p=p$
\item[$\bullet$] $V(X)= P(X=0)\times (0-E(X))^2+P(X=1)\times (1-E(X))^2=(1-p)\times (0-p)^2+p\times (1-p)^2=p^2(1-p)+p(1-p)^2=p(1-p)(p+1-p)=p(1-p)$
\end{itemize}
\end{demo}
\begin{comment}
\begin{rem}[Algorithmique]
Algorithme de simulation d'une épreuve de Bernoulli de paramètre $p$.

\begin{algorithm}
\Donnees{$p$ : nombre décimal entre 0 et 1 ;}
\Traitement{
$t$ prend une valeur aléatoire décimale entre 0 inclus et 1 exclu ;\\
\Si{$t<p$}{Renvoyer "Succès" ; }
\Sinon{Renvoyer "\'Echec" ;}
}
\end{algorithm}
\end{rem}

\begin{rem}{Exemple de programmation en langage python}
\begin{lstlisting}[language=python,frame=single,keywordstyle=\color{blue},numberstyle=\color{green},stringstyle=\color{orange},identifierstyle=\color{purple}]
import random
def bernoulli(p):
	t=random.random()
	if t<p:
		return 1
	else:
		return 0
\end{lstlisting}
\end{rem}
\end{comment}

\section{Schéma de Bernoulli}

\begin{defn}
Deux expériences sont dites indépendantes si le résultat de l'une n'influe pas sur le résultat de l'autre.
\end{defn}

\begin{exemple}
il y a indépendance lorsqu'on lance deux fois de suite une pièce de monnaie.
\end{exemple}

\begin{defn}
On considère une expérience aléatoire ne comportant que deux issues ; l'une appelée \og{}succès\fg{} et l'autre appelée \og{}échec\fg{}. On note $p$ la probabilité de succès et $q$ la probabilité d'échec ($q=1-p$). La répétition de cette expérience $n$ fois de manière indépendante constitue \emph{un schéma de Bernoulli} de paramètres $n$ et $p$.
\end{defn}



\begin{prop}
Dans un schéma de Bernoulli, la probabilité d'une liste de résultats est le produit des probabilités de chaque résultat
\end{prop}


\begin{exemple}
%\begin{center}
 %\includegraphics[width=6cm]{arbre.png}
%\end{center}
On contrôle la qualité d'un produit sur une chaîne de production. On prélève 3 produits au hasard. On suppose que les prélèvements sont indépendants. Statistiquement, chaque produit a une probabilité $p=0,05$ d'être défectueux.\\
Sur l'arbre ci-dessus représentant un schéma de Bernoulli de paramètres $n=3$ et $p=0,05$, la probabilité d'avoir les deux premières expériences qui donnent un succès et la dernière qui donne un échec est $P(SS\bar{S})=0,05^2\times 0,95\approx 0,002$\\
soit 0,2 \% de chances d'avoir deux produits défectueux sur les trois prélevés.
\end{exemple}

\section{Loi binomiale}


\begin{defn}
Soient $n$ et $k$ deux entiers naturels avec $k\leq n$. On note $\C{n}{k}$ (on dit \og{}$k$ sous $n$\fg{}) le nombre de manières d'obtenir $k$ succès et $n-k$ échecs pour $n$ répétitions indépendantes de la même expérience de Bernoulli. 
\end{defn}




\begin{exemple}
%\begin{center}
 %\includegraphics[width=8cm]{arbre.png}
%\end{center}

\noindent $\C{3}{3}=1$ : il y a une seule manière d'obtenir 3 succès lors de la répétition de 3 épreuves identiques indépendantes.\\
$\C{3}{2}=3$ : il y a trois manières d'obtenir 2 succès et un échec lors de la répétition de 3 épreuves identiques indépendantes ($SSE$ ; $SES$ ; $ESS$).
\end{exemple}

\begin{defn}
On considère une épreuve de Bernoulli de paramètre $p\in [0;1]$. On répète $n$ fois ($n\geq 1$) cette expérience indépendamment et on note $X$ la variable aléatoire qui compte le nombre de succés. On dit alors que la variable aléatoire $X$ suit une loi \emph{binomiale} de paramètres $n$ et $p$ et on note $X\sim \mathcal{B}(n;p)$.
\end{defn}



\begin{prop}
 Si $X$ est une variable aléatoire qui suit une loi binomiale de paramètres $n$ et $p$, alors la probabilité d'obtenir $k$ succès avec $k\in\{0;1;2;...;n\}$ est $P(X=k)=\C{n}{k}p^k(1-p)^{n-k}$.
\end{prop}


\begin{demo}
Il y a $\C{n}{k}$ manières d'obtenir $k$ succès dans $n$ répétitions d'expériences identiques et indépendantes. La probabilité de chacune de ces événements qui sont évidemment incompatibles est $p^k(1-p)^{n-k}$ d'où le résultat.
\end{demo} 

\begin{exemple}
On considère le problème précédent de test des produits d'une chaîne de production. Les prélèvements étant supposés indépendants les uns des autres, l'expérience constitue un schéma de Bernoulli de paramètres $n=3$ et $p=0,05$. La variable aléatoire $X$ qui compte le nombre de succès suit la loi binomiale de paramètres $n=3$ et $p=0,05$.\\
On a $P(X=2)=P(SS\bar{S}\cap S\bar{S}S\cap \bar{S}SS)$\\
 car trois chemins permettent d'obtenir deux succès c'est à dire deux objets défectueux.\\
D'où $P(X=2)=P(SS\bar{S})+P(S\bar{S}S)+P(\bar{S}SS)$\\
donc $P(X=2)=0,05^2\times 0,95+0,95\times 0,05\times 0,95+0,95\times 0,05^2=3\times 0,05^2\times 0,95=0,007$\\
soit une probabilité très faible de 0,007 d'avoir deux produits défectueux.
\end{exemple}

\begin{rem}{Calcul pratique de $P(X=k)$ et $P(X\leq k)$}
Soit $X$ une variable aléatoire de paramètres $n$ et $p$. Pour $k$ allant de 0 à $n$, pour calculer $P(X=k)$ ou $P(X\leq k)$, on utilise une calculatrice :
\begin{itemize}
\item[$\bullet$] Sur Texas instrument : aller dans le menu \fbox{2nd}\fbox{distrib}, choisir \fbox{binomFdp} et taper \fbox{n}\fbox{,}\fbox{p}\fbox{,}\fbox{k}\fbox{)} pour calculer $P(X=k)$ et choisir \fbox{binomFRép} et taper \fbox{n}\fbox{,}\fbox{p}\fbox{,}\fbox{k}\fbox{)} pour calculer $P(X\leq k)$.
\item[$\bullet$] Sur Casio : aller dans le menu \fbox{STAT} puis \fbox{DIST} puis \fbox{BINM}. Sélectionner alors \fbox{Bpd} puis \fbox{Var} pour variable, puis entrer alors $k$  dans la ligne \og{}x\fg{}, $n$ dans la ligne \og{}numtrial\fg{} et $p$ dans la ligne \og{}p\fg{} puis aller sur \og{}execute\fg{} pour valider et calculer ainsi $P(X=k)$. Pour le calcul de  $P(X\leq k)$, on utilisera \fbox{Bcd} au lieu de \fbox{Bpd}.
\end{itemize}
\end{rem}

\begin{exemple}
On considère une variable aléatoire $X$ suivant la loi binomiale de paramètres $n=4$ et $p=0,4$.\\
Sur TI, la probabilité $P(X\leq 3)$ est donnée par \texttt{binomFrép(4,0.4,3)}.
\end{exemple}

\begin{rem}
\begin{itemize}
\item[$\bullet$] On a $P(X<3)=P(X\leq 2)$\\
\item[$\bullet$] pour calculer $P(X>k)$, on calcule $1-P(X\leq k)$.
\end{itemize}
\end{rem}

\begin{prop}
Soit $X$ une variable aléatoire qui suit la loi binomiale de paramètres $n$ et $p$ :
\begin{itemize}
\item[$\bullet$] $E(X)=np$ ;
\item[$\bullet$] $V(X)=np(1-p)$ et $\sigma(X)=\sqrt{np(1-p)}$.
\end{itemize}
\end{prop}

\begin{demo}
Pour rappel si $X_1,X_2,\cdots X_n$ sont des variables aléatoires indépendantes qui suivent une loi de Bernouilli de paramètre $p$ alors $X_1+X_2+\cdots + X_n$ est une variable aléatoire qui suit une loi Binomiale de paramètre $\calb(n,p)$.

Ainsi soit $Z$ une variable aléatoire qui suit une loi Binomiale de paramètre $\calb(n,p)$ alors $E(Z)=E(X_1+X_2+\cdots + X_n)=E(X_1)+E(X_2)+\cdots +E(X_n)=p+p+\cdots+p=np$.

De même pour la variance $V(Z)=V(X_1+X_2+\cdots + X_n)=V(X_1)+V(X_2)+\cdots +V(X_n)=p(1-p)+p(1-p)+\cdots+p(1-p)=np(1-p)$.

On peut donc déduire aussi l'écart-type $\sigma(Z)=\sqrt{np(1-p)}$.
\end{demo}


\begin{exemple}
Pour le problème de la chaîne de production, en prélevant $n=100$ produits indépendamment, la loi binomiale a pour paramètres $n=100$ et $p=0,05$. On a alors $E(X)=np=100\times 0,05=5$ ce qui signifie que l'on peut prévoir 5 produits défectueux pour un prélèvement de 100 produits indépendants.
\end{exemple}

\begin{comment}
\begin{rem}
Algorithme de simulation d'une loi binomiale de paramètres $n$ et $p$.

\begin{algorithm}
\Donnees{$p$ : nombre décimal entre 0 et 1 ; $n$ : entier naturel ;}
\Traitement{
$c\leftarrow 0$;\\
\Pour{$k$ de 1 jusque $n$}
  {$t$ prend une valeur aléatoire décimale entre 0 inclus et 1 exclu;\\
  \Si{$t<p$}{$c\leftarrow c+1$;\\}
  }
}
\end{algorithm}
\end{rem}
\end{comment}

\section{Coefficients binomiaux}


\begin{defn}
On appelle \emph{coefficients binomiaux} l'ensemble des nombres $\C{k}{n}$ pour $n\in\mathbb{N}$ et $k\in\mathbb{N}$.
\end{defn}

\begin{rem}[Utilisation de la calculatrice]
\begin{itemize}
\item[Casio : ] Touche \fbox{nCr}. Taper sur \fbox{OPTN} puis \fbox{PROB} puis \fbox{nCr}.\\
Syntaxe : 7\fbox{nCr}4
\item[TI : ] Touche \fbox{Combinaison}. Taper sur \fbox{math} puis aller dans le menu \fbox{PROB} puis choisir \fbox{Combinaison}.\\
Syntaxe : 7\fbox{Combinaison}4
\item[Numworks : ] menu boîte à outils puis \fbox{Dénombrement} puis \fbox{binomial(n,k)} (Attention : nom mal choisi, cela n'a rien à voir avec une loi binomiale).\\

\end{itemize}
\end{rem} 


\begin{prop}
\begin{itemize}
\item[$\bullet$] $\C{n}{k}=\C{n}{n-k}$ ;
\item[$\bullet$] $\C{n}{k}+\C{n}{k+1}=\C{n+1}{k+1}$
\end{itemize}
\end{prop}




\section{Triangle de Pascal}
 Le triangle de Pascal, du nom de Blaise Pascal, mathématicien français du XVII\up{e} siècle qui le\og{}redécouvra\fg{} en Occident (car il était connu avant en Orient et au Moyen-Orient) est une disposition permettant de visualiser et de calculer les coefficients binomiaux et qui s'appuie sur la formule précédente.
\begin{center}
 \begin{tabular}{|p{1.5cm}|p{1.5cm}|p{1.5cm}|p{1.5cm}|p{1.5cm}|p{1.5cm}|}
\hline
  $\C{0}{0}=1$ & & & & &\\
\hline
$\C{1}{0}=1$ & $\C{1}{1}=1$ & & & &\\
\hline
$\C{2}{0}=1$ & $\C{2}{1}=2$ & 1 & & &\\
\hline
$\C{3}{0}=1$ & $\C{3}{1}=3$ & 3 & 1 & & \\
\hline
1 & 4 & 6 & 4 & 1 &\\
\hline
1 & 5 & 10 & 10 & 5 & 1\\
\hline
 \end{tabular}
\end{center}
Explication de la construction : le nombre de la ligne $n$ et de la colonne $k$ est le coefficient binomial $\C{n-1}{k-1}$. Il est obtenu en ajoutant le nombre situé au dessus (ligne $n-1$ et colonne $k$) au nombre de la colonne et de la ligne précédente (ligne $n-1$ et colonne $k-1$).\\
Par exemple, $\C{3}{1}=3$ est la somme de $\C{2}{1}=2$ et de $\C{2}{0}=1$.
\begin{center}
 \begin{tabular}{|c|c|c|c|c|c|}
\hline
  1 & & & & &\\
\hline
1 & 1 & & & &\\
\hline
1 & 2 & 1 & & &\\
\hline
1 & 3 & 3 & 1 & & \\
\hline
1 & 4 & 6 & 4 & 1 &\\
\hline
1 & 5 & 10 & 10 & 5 & 1\\
\hline
 \end{tabular}
\end{center}

\newpage

\begin{comment}
\begin{sanscadre}{Exemple de programmation du triangle de Pascal en langage python}
\begin{lstlisting}[language=python,frame=single,keywordstyle=\color{blue},numberstyle=\color{green},stringstyle=\color{orange},identifierstyle=\color{purple}]
def trianglePascal(n):
	L=[1]
	L=L+[0 for k in range(n)]
	t=[L]
	for k in range(1,n+1):
		L=[1]
		for l in range(1,k+1):
			L.append(t[k-1][l]+t[k-1][l-1])
		L=L+[0 for l in range(k,n)]
		t=t+[L]
	return t
\end{lstlisting}
\end{sanscadre}
\end{comment}


\section{Loi géométrique}


\begin{defn}
Soit $E=\mathbb{N}^{*}$. Soit $p\in [0;1]$. On dit qu'une variable aléatoire $X$ définie sur $E$ suit une loi géométrique de paramètre $p$ lorsqu'elle donne le nombre d'expériences de Bernoulli de paramètre $p$ répétées indépendamment nécessaires pour obtenir un premier succès. On note alors $X\sim \mathcal{G}(p)$.
\end{defn}


\begin{prop}
Si $X$ suit une loi géométrique de paramètre $p$ alors :
\begin{itemize}
\item[$\bullet$] pour tout $k\in\mathbb{N}^{*}$, 
$$P(X=k)=p(1-p)^{k-1}$$
\item[$\bullet$] $E(X)=\frac{1}{p}$.
\end{itemize}
\end{prop}


\begin{prop}
$X$ suit une loi géométrique si et seulement si pour tout $m,n\in\mathbb{N}^{*}$, $P_{X>m}(X>m+n)=P(X>n)$.\\
Cette propriété illustre que la loi géométrique est une loi dite \og{}sans mémoire\fg{}. Cette propriété montre aussi que les seules lois discrètes sans mémoire sont les lois géométriques.
\end{prop}


\begin{demo}[Preuve de la condition nécessaire]
$P_{X>m}(X>m+n)=\frac{P(X>m+n\cap X>m)}{X>m}=\frac{P(X>m+n)}{X>m}$ car $X>m$ est nécessairement réalisé si $X>m+n$.\\
Remarquons ensuite que $P(X>m)=1-P(X\leq m)=1-\sum_{k=1}^{m}p(1-p)^{k-1}=1-p\sum_{k=0}^{m-1}(1-p)^k=1-p\frac{1-(1-p)^m}{1-(1-p)}=1-p\frac{1-(1-p)^m}{p}=1-(1-(1-p)^m)=(1-p)^m$\\
D'où $P_{X>m}(X>m+n)=\frac{P(X>m+n)}{X>m}=\frac{(1-p)^{m+n}}{(1-p)^m}=(1-p)^n$.\\
Or $P(X>n)=(1-p)^n$ d'où l'égalité.
\end{demo}

\end{document}
